\documentclass[aspectratio=169]{beamer}

% -------------------- Theme --------------------
\usetheme{Madrid}
\usecolortheme{default}
\setbeamertemplate{navigation symbols}{}
\setbeamertemplate{footline}[frame number]

% -------------------- Packages --------------------
\usepackage{amsmath,amssymb,bm}
\usepackage{booktabs}
\usepackage{tikz}
\usetikzlibrary{arrows.meta,positioning,calc,decorations.pathreplacing}

% -------------------- TikZ styles --------------------
\tikzset{
  var/.style={circle, draw, thick, minimum size=8mm, align=center},
  obs/.style={circle, draw, thick, fill=blue!8, minimum size=8mm, align=center},
  treat/.style={circle, draw, thick, fill=blue!15, minimum size=8mm, align=center},
  out/.style={circle, draw, thick, fill=orange!20, minimum size=8mm, align=center},
  conf/.style={circle, draw, thick, fill=purple!12, minimum size=8mm, align=center},
  hidden/.style={circle, draw, dashed, thick, fill=gray!10, minimum size=8mm, align=center},
  arrow/.style={-Latex, thick},
  badpath/.style={-Latex, thick, red!70},
  note/.style={rectangle, rounded corners, draw, fill=yellow!12, align=left}
}

% -------------------- Macros --------------------
\newcommand{\E}{\mathbb{E}}
\newcommand{\Prob}{\mathbb{P}}
\newcommand{\indep}{\perp\!\!\!\perp}
\newcommand{\ATE}{\mathrm{ATE}}
\newcommand{\ATT}{\mathrm{ATT}}
\newcommand{\doop}{\mathrm{do}}

% -------------------- Metadata --------------------
\title[Causal Inference]{Lecture Title Here}
\subtitle{Causal Inference in the AI Era}
\author{Dan Vilenchik}
\institute{Ben-Gurion University of the Negev}
\date{\today}

\begin{document}

\begin{frame}
  \titlepage
\end{frame}

\begin{frame}{Today’s goal}
\begin{itemize}
  \item Intuition first.
  \item Formal estimand second.
  \item Assumptions and identification third.
  \item Estimator last.
\end{itemize}
\end{frame}

\begin{frame}{A simple causal question}
\begin{center}
\begin{tikzpicture}[node distance=2.2cm]
  \node[treat] (T) {$T$};
  \node[out, right=of T] (Y) {$Y$};
  \draw[arrow] (T) -- (Y);
\end{tikzpicture}
\end{center}

\[
\text{Goal: } \E[Y(1)-Y(0)]
\]
\end{frame}

\begin{frame}{Observed vs. counterfactual}
\[
Y^{obs}=T Y(1)+(1-T)Y(0)
\]

\begin{itemize}
  \item If $T=1$, we observe $Y(1)$ and miss $Y(0)$.
  \item If $T=0$, we observe $Y(0)$ and miss $Y(1)$.
\end{itemize}
\end{frame}

\begin{frame}{Common mistake}
\begin{block}{Do not confuse identification and estimation}
Identification says an observable formula equals a causal estimand.\\
Estimation says how we compute that formula from finite data.
\end{block}
\end{frame}

\begin{frame}{Recap}
\begin{itemize}
  \item What was the causal question?
  \item What was the estimand?
  \item What assumption identified it?
  \item What estimator would we use?
\end{itemize}
\end{frame}

\end{document}
